% !TeX root = ../main.tex
% Seismic and dynamic design. Moved to Part IV: seismic action is an
% imposed load case handled by design and analysis, not a property of
% the ground.
\chapter{Seismic and Dynamic Design}
\label{ch:seismic}

Seismic action is distinguished from the other loads considered in this book
by the manner in which it is applied. The ground imposes a deformation, and
the structure either accommodates that deformation or fails to; it is not
loaded in the sense that a lining is loaded by ground pressure. The
consequence, developed through the chapter, is that increased strength and
stiffness are frequently the wrong response, and that the governing parameter
is the relative stiffness of structure and ground rather than the capacity of
the structure alone. Blast vibration during construction, the other dynamic
action a tunnel commonly imposes on its surroundings, is governed by control
of the source rather than by imposed ground deformation, and is treated with
the excavation methods in Section~\ref{sec:vibration-environment}.

\section{Seismic demand and analysis}
\label{sec:seismic-demand}

\subsection{Performance of underground structures}
\label{sec:seismic-performance}

Underground structures perform well in earthquakes, and the reason is
kinematic rather than structural. A building responds dynamically, amplifying
ground motion through its own inertia; a tunnel is confined by ground on all
sides, its mass small relative to the surrounding soil, so it moves with the
ground rather than responding to it. Compilations of earthquake damage to
underground structures~\citep{hashash2001seismic} confirm this consistently:
deep bored tunnels in competent rock come through strong shaking
with little or no damage.

The exceptions define the design problem. Damage concentrates where cover is
shallow, at portals, at junctions between structures of differing stiffness, in
liquefiable ground, and above all where a fault ruptures through the
alignment. Older tunnels with unreinforced masonry or brittle linings are more
vulnerable than modern reinforced ones, and tunnels already deteriorated by
water ingress and corrosion are more vulnerable again, an interaction
between seismic vulnerability and asset condition that
Chapter~\ref{ch:sustainability} identifies as a resilience issue.

The main distinction is between ground shaking and ground failure.
Shaking imposes transient deformation as waves pass; the structure deforms,
develops strain, and returns. Ground failure (fault rupture, liquefaction,
landsliding) imposes permanent displacement that the structure must
accommodate or fail. These require different treatments and have to be checked
separately, not combined into a single seismic check. They are coupled in one direction only: ground failure
can remove the confinement on which the shaking argument depends, so a site
exposed to both requires the shaking check to be repeated for the degraded
condition.

\subsection{Hazard and deformation demand}
\label{sec:hazard-deformation}

Seismic hazard assessment produces a design ground motion for a stated return
period, and underground structures are usually designed for two levels: an
operating earthquake, more frequent, after which the asset remains in service;
and a maximum design earthquake, rare, after which life safety is assured and
repairable damage is accepted. Intraplate regions of low to moderate
seismicity present a particular framing problem, since the hazard is dominated
by rare large events rather than by frequent moderate ones. The operating
earthquake then becomes so small that it imposes no effective check, the two-level
framework reduces in practice to a single rare-event case, and because the
design motion grows steeply with return period the result is unusually
sensitive to the return period selected.

The demand on the structure is deformation, not force. Three modes are
distinguished (Figure~\ref{fig:seismic-modes}). \emph{Axial} extension and compression occur as a wave travels
along the tunnel axis, straining the lining longitudinally.
\emph{Longitudinal bending} occurs as the wave curves the alignment.
\emph{Ovaling} of a circular section, or \emph{racking} of a rectangular one,
occurs as a shear wave propagates perpendicular to the axis, distorting the
cross-section.

\begin{figure}[htbp]
\centering
\datafigure[\textwidth]{ch20/seismic-modes}
\caption{Deformation imposed on a tunnel by seismic waves (schematic, exaggerated): axial extension and compression in plan (with the axial strain along the tunnel, carmine), curvature of the alignment, ovaling of a circular section and racking of a rectangular one under shear. Dashed outlines are undeformed; arrows give the direction of wave travel or of shear.}
\label{fig:seismic-modes}
\end{figure}

Free-field deformation methods take the strain that would occur in the ground
without the tunnel and impose it on the structure. The wave-passage approach of
\citet{newmark1968problems} gives axial and curvature strain from the peak particle velocity
(the measure of ground vibration also used for blasting,
Section~\ref{sec:vibration-environment}), the peak particle
acceleration, and an assumed wave speed and angle of incidence. For a shear
wave of peak particle velocity $v_{\max}$ and propagation velocity $C_s$
incident at an angle $\theta$ to the tunnel axis, the free-field axial
strain is
\[
\varepsilon_a = \frac{v_{\max}}{C_s}\,\sin\theta\cos\theta ,
\]
which is greatest at $45^\circ$, where it equals $v_{\max}/(2C_s)$. The angle
is unknown, so the design case takes that worst orientation, which is
conservative and costs little in analysis, since the envelope follows from a
single closed-form evaluation over the range of incidence angles. A particle
velocity of 0.3~m/s in ground with
$C_s = 1000$~m/s gives a strain of at most $1.5 \times 10^{-4}$
(Figure~\ref{fig:free-field-strain}). Imposing the
free-field strain is appropriate where the lining is flexible relative to the
ground and the ground is uniform over a wavelength. At stiffness transitions,
portals and junctions the assumption errs on the unconservative side, because
the deformation concentrates at the interface rather than spreading with the
free field, and it has to be replaced by an interaction analysis; these are
the critical elements of Section~\ref{sec:seismic-structures-criteria}.

\begin{figure}[htbp]
\centering
\datafigure[\textwidth]{ch20/free-field-strain}
\caption{Free-field axial strain from a shear wave after \citet{newmark1968problems}: against the angle of incidence for $v_{\max} = 0.3$~m/s, and at the worst angle, $v_{\max}/(2C_s)$, against particle velocity, each for three values of $C_s$. The dot is the example of the text. The free-field shear strain that drives ovaling, $v_{\max}/C_s$, is twice this peak.}
\label{fig:free-field-strain}
\end{figure}

\subsection{Simplified analysis methods}
\label{sec:simplified-seismic}

Ovaling and racking are handled by closed-form
solutions~\citep{wang1993seismic,penzien2000seismically} that are governed
by a single ratio (Figure~\ref{fig:ovaling-flexibility}). The \emph{flexibility ratio} compares the shear stiffness
of the ground to the bending stiffness of the lining, and it determines
whether the lining follows the ground or resists it. At high flexibility ratio
the lining conforms: its distortion exceeds the free-field distortion and
approaches that of an unlined opening, $4(1-\nu)$ times the free field, while
the induced forces are small. The lining distorts as the free field only near
a flexibility ratio of about 7 to 9. At low flexibility ratio the lining
resists, distorts less, and attracts large moments. The \emph{compressibility ratio} has the
same function for the axial mode (Section~\ref{sec:charts-and-lining-solutions}).

\begin{figure}[htbp]
\centering
\datafigure[\textwidth]{ch20/ovaling-flexibility}
\caption{Full-slip ovaling of a circular lining after \citet{wang1993seismic}, against the flexibility ratio $F$ of Appendix~\ref{app:closed-form} (Wang's ratio is $F(1-\nu)/6$). The lining distorts as the free field near $F \approx 7$ to 9, less when stiffer, and up to $4(1-\nu)$ times as much, the distortion of a cavity, when flexible; its moment, $K_1 E R^2 \gamma_{\max}/[6(1+\nu)]$, falls steadily with $F$.}
\label{fig:ovaling-flexibility}
\end{figure}

The design implication mirrors that of static lining design in
Chapter~\ref{ch:linings}: a flexible lining is generally the better seismic
structure. Stiffening a lining to resist seismic distortion increases the
forces it must carry, and the more effective response is usually to allow the
distortion and detail for it.

Interface assumptions bracket the answer. Full-slip between lining and ground
gives lower thrust and higher moment; no-slip gives the reverse. Actual behaviour lies between, and both should be checked, since the governing
condition differs between reinforcement design and section capacity.

\section{Seismic design and permanent displacement}
\label{sec:seismic-design-displacement}

\subsection{Dynamic numerical analysis}
\label{sec:dynamic-analysis}

Where the simplified methods do not apply (irregular geometry, strongly
layered ground, soil-structure interaction at a portal or a station box),
dynamic numerical analysis is required. Three inputs dominate the result.

Input motion must be defined at the model base, which requires deconvolving a
surface motion downward through the soil profile; errors here propagate
directly. Boundary conditions must absorb outgoing waves rather than reflect
them, since a reflecting boundary traps energy in the model and produces
spurious amplification. Damping must represent energy dissipation without
suppressing the response. The formulation in most common use, Rayleigh
damping, is proportional to mass and stiffness, so that the damping ratio it
produces at circular frequency $\omega$ is $a_R/(2\omega) + b_R\,\omega/2$. It
matches a target ratio at only two frequencies
(Figure~\ref{fig:rayleigh-damping}), falls below the target
between them and rises above it outside them, and the two frequencies have
therefore to be chosen to bracket the range that matters to the structure and
stated with the result.

\begin{figure}[htbp]
\centering
\datafigure[0.62\textwidth]{ch20/rayleigh-damping}
\caption{Rayleigh damping set to a target of 5\,\% at 1 and 8~Hz (illustrative frequencies). The ratio falls to about 3\,\% between them and rises steeply outside them, through its mass part at low frequency and its stiffness part at high.}
\label{fig:rayleigh-damping}
\end{figure}

An analysis that has not been checked against a simplified solution for a
comparable case should not be relied on, for the reasons set out in
Chapter~\ref{ch:numerical-modelling}. Here the check has a specific purpose:
running the model first on a reduced configuration for which the closed form
holds, uniform ground and a circular section under free-field shear, tests
the input motion, the absorbing boundaries and the damping before the
irregular geometry that motivated the model is introduced.

\subsection{Permanent ground displacement}
\label{sec:permanent-displacement}

Fault rupture through an alignment imposes displacement that no lining can
resist. The design objective becomes accommodation and repairability rather
than protection. Options include oversizing the section through the fault zone
so that the operating envelope survives the displacement; articulating the
lining with flexible joints that concentrate the deformation where it can be
tolerated; using a ductile lining that deforms without collapse; and accepting
damage in a defined length while ensuring it can be repaired
(Figure~\ref{fig:fault-crossing}).

Two quantities decide among these options: the expected displacement and the
uncertainty in where the rupture will cross the alignment. Oversizing works
only while the displacement remains small relative to the clearance
provided, and the enlarged section has to span the whole width within which
the rupture may cross. Articulation fails if the rupture falls between the
designed joints, so it suits a narrow fault zone whose location is well
constrained. A ductile lining provides deformation capacity rather than
displacement capacity: it spreads a modest offset without collapse but cannot
absorb a large one. Accepting and repairing damage is the option left where
the displacement is large or its location poorly known.

\begin{figure}[htbp]
\centering
\datafigure[\textwidth]{ch20/fault-crossing}
\caption{Four ways of accommodating fault displacement, in long section across a fault zone (grey) with the rupture (line), schematic: an enlarged section within which the operating envelope (dashed) can be realigned; short rings with flexible joints; a ductile lining spreading the offset; and a defined length accepted as damaged and repaired.}
\label{fig:fault-crossing}
\end{figure}

Both quantities come from the geological assessment of fault activity, slip
rate and the location and width of the fault zone, a specialist input that must be obtained early,
since it determines the alignment as much as the design.

Liquefaction affects tunnels differently from buildings. A buoyant structure
in liquefied soil floats: a shallow tunnel whose overall unit weight, structure
and enclosed air together, is well below that of the liquefied soil can uplift
substantially. Loss of lateral support during
liquefaction also removes the confinement on which the flexible-lining
argument depends. Mitigation is by ground improvement to prevent liquefaction,
or by increasing the structure's weight, and the assessment must consider the
tunnel empty as well as in service.

\subsection{Structures and performance criteria}
\label{sec:seismic-structures-criteria}

The bored tunnel is rarely the critical element. Portals combine a slope
problem with a stiffness transition; station boxes are stiff rectangular
structures that rack rather than ovalise, and because a frame cannot shed
shear distortion into a change of ring shape they attract substantially
higher moments than a circular section in the same ground; shafts impose a vertical stiffness discontinuity; and
junctions between structures of differing stiffness concentrate deformation at
the interface. Seismic design effort should be directed at these rather than
distributed uniformly along the alignment.

Performance-based criteria make the objective explicit. An operational level
requires the asset to remain in service after the event. Life safety requires
that people can evacuate and that collapse does not occur, accepting damage
and closure. Matching performance level to hazard level and to asset
criticality (operational for the operating earthquake, life safety
for the maximum design earthquake) is the standard framing, and it connects
to the functionality-over-time treatment in
Chapter~\ref{ch:sustainability}.

This chapter treated seismic action as a deformation imposed by the ground.
The free-field strain follows from the particle velocity and the wave speed,
and the flexibility ratio decides how much of the ground distortion the
lining takes up and what moment it attracts. A flexible, ductile lining is
therefore generally the better seismic structure, and the design effort
belongs at portals, station boxes, shafts, junctions and fault crossings,
where a stiffness contrast or a permanent displacement governs. The inputs
to these checks, the design motion, the wave speed and the fault
displacement, are uncertain, and Chapter~\ref{ch:integrated-design} deals
with how such uncertainty is carried through a design.

\chapterreferences
